% A \write's TEXT IS EXPANDED IN NO MODE, and \tracingcommands says so.
%
% Every command traced while the text of a \write is being expanded carries
% the mode `no mode', which is none of the four:
%
%   \shipout\hbox{\write16{\number\count5}}   {no mode: \number}
%   \immediate\write16{\number\count5}        {no mode: \number}
%
% Both kinds: the deferred one is expanded when the page is shipped and the
% immediate one where it stands, and both are in no mode. Because the mode
% is named only where it has changed, this also makes the NEXT trace line
% after a \write name its mode again -- which is how the difference shows
% up in trip far more often than the \write lines themselves do.
%
% NO OTHER EXPANDED TEXT DOES THIS. Measured, all in vertical mode and all
% drawing the mode that stands: \message, \edef, \special, \mark, and a
% token parameter's text (\errhelp).
%
% \immediate HIDES THE COMMAND IT REACHES OVER. It reads that command
% itself, so the main loop never obeys it and never traces it:
%
%   \immediate\write16{x}      {\immediate}
%   \immediate\openout15=f     {\immediate}
%   \immediate\closeout15      {\immediate}
%   \immediate\special{x}      {\immediate} {\special}
%   \immediate\relax           {\immediate} {\relax}
%
% -- the last two because \special and \relax are not among the three
% \immediate accepts, so they are put back and met again by the main loop.
\tracingonline=1 \count5=3 \def\q{Q}
\tracingcommands=2
\shipout\hbox{\write16{[a\number\count5\q]}}
\immediate\write16{[b\number\count5]}
\message{[m\number\count5]}
\edef\z{[e\number\count5]}
\special{[s\number\count5]}
\mark{[k\number\count5]}
\errhelp{[h\number\count5]}
\immediate\relax
\immediate\openout15=rfout
\immediate\closeout15
\immediate\special{x}
\tracingcommands=0
\end
